\providecommand{\atlasref}[2]{#2}
\section{Całka Lebesgue'a}\label{ch:calkowanie-lebesgue}

Ograniczamy się do podstawowych twierdzeń o zbieżności całek.
Konstrukcję całki, mierzalność, całkowalność i rachunek całek
omawia skrypt \emph{Teoria miary}, w szczególności
\atlasref{ME-06::thm:simple-up}{ME-06: aproksymacja funkcji mierzalnych},
\atlasref{ME-07::ch:calkowanie-lebesgue}{ME-07: konstrukcja całki Lebesgue'a}
oraz \atlasref{ME-08::thm:mct}{ME-08: twierdzenia graniczne}.
Dalszą teorię rozwinięto w rozdziałach ME-09--ME-14
(rodzaje zbieżności, regularność, całki z parametrem,
miary produktowe, zamiana zmiennych i przestrzenie $L^p$)
oraz ME-25 (różniczkowanie całek i miar).
W szczególności zob.
\atlasref{ME-11::ch:parametry-calki}{ME-11: całki zależne od parametru},
\atlasref{ME-12::thm:fubini-12}{ME-12: twierdzenia Tonellego i Fubiniego},
\atlasref{ME-13::thm:change-variables-13}{ME-13: zamiana zmiennych},
\atlasref{ME-14::thm:riesz-fischer-14}{ME-14: przestrzenie $L^p$}
i \atlasref{ME-25::thm:lebesgue-diff-25}{ME-25: różniczkowanie całek}.
Nie powtarzamy konstrukcji ani dowodów tych wyników.

Niech $(X,\mathcal F,\mu)$ będzie dowolną przestrzenią z miarą.
Rozważamy funkcje mierzalne; $L^1(\mu)$ oznacza przestrzeń
klas funkcji całkowalnych modulo równość prawie wszędzie.
Nie zakładamy zupełności miary ani $\sigma$-skończoności.

\subsection{Dlaczego nie wolno bezwarunkowo przestawiać granicy i całki?}\label{ch:twierdzenia-graniczne}
Na $(0,1)$ z miarą Lebesgue'a $\lambda_1$ funkcje
$f_n=n\chi_{(0,1/n)}$ zbiegają punktowo do zera:
dla ustalonego $x>0$ i $n>1/x$ mamy $f_n(x)=0$.
Jednak $\int_{(0,1)} f_n\,\mathrm d\lambda_1
=n\lambda_1((0,1/n))=1$ dla każdego $n$, podczas gdy
całka funkcji granicznej wynosi zero.
Pokazuje to zarazem, że w lemacie Fatou nierówność
może być ostra.
Poniższe twierdzenia określają warunki zamiany granicy
i całkowania; pełna teoria znajduje się w
\atlasref{ME-08::thm:mct}{ME-08}. Symbol $f_n\uparrow f$
oznacza niemalejącą zbieżność punktową.

\subsection{Twierdzenie Beppo Leviego}


\begin{twierdzenie}[O zbieżności monotonicznej; Beppo Levi]
\label{thm:mct}
Niech $(f_n)_{n\ge1}$ będzie ciągiem mierzalnych funkcji
$f_n:X\to[0,\infty]$, takim że dla każdego $x\in X$
$f_n(x)\le f_{n+1}(x)$ dla każdego $n\ge1$. Niech
$f(x)=\lim_{n\to\infty}f_n(x)=\sup_nf_n(x)$.
Wtedy $f$ jest mierzalna i
\begin{equation}\label{eq:mct}
\int_Xf(x)\,\mathrm d\mu(x)
=\lim_{n\to\infty}\int_Xf_n(x)\,\mathrm d\mu(x)
=\sup_{n\ge1}\int_Xf_n(x)\,\mathrm d\mu(x).
\end{equation}
Dopuszczamy wartość $+\infty$ i nie wymagamy
żadnego założenia o $\mu(X)$.
\end{twierdzenie}
\noindent\emph{Dowód:} zob. \atlasref{ME-08::thm:mct}{Teoria miary, rozdział 8 (dowód tego twierdzenia)}.

\begin{wniosek}[Konstruktywne obliczanie całki]
\label{cor:simple-approx-integral}
Niech $(X,\mathcal F,\mu)$ będzie przestrzenią z miarą.
Jeżeli $f:X\to[0,\infty]$ jest mierzalna, a
$(s_n)$ jest \emph{dowolnym} ciągiem nieujemnych
funkcji prostych mierzalnych takim, że $s_n\uparrow f$
wszędzie, to
$\int_Xf\,\mathrm d\mu(x)=
\lim_{n\to\infty}\int_Xs_n\,\mathrm d\mu(x)$.
W szczególności można użyć funkcji z
\atlasref{ME-06::thm:simple-up}{ME-06: twierdzeniu o aproksymacji funkcjami prostymi}.
\end{wniosek}
\noindent\emph{Dowód:} zob. \atlasref{ME-08::cor:simple-approx-integral}{Teoria miary, rozdział 8 (dowód tego twierdzenia)}.

\begin{twierdzenie}[Zbieżność monotoniczna prawie wszędzie]
\label{thm:mct-ae}
Niech $f_n,f:X\to[0,\infty]$ będą mierzalne.
Jeśli dla każdego $n\ge1$ zachodzi $f_n\le f_{n+1}$ prawie wszędzie oraz $f_n\to f$
$\mu$-prawie wszędzie, to zachodzi równość
\eqref{eq:mct}.
\end{twierdzenie}
\noindent\emph{Dowód:} zob. \atlasref{ME-08::thm:mct-ae}{Teoria miary, rozdział 8 (dowód tego twierdzenia)}.

\subsection{Lemat Fatou}


\begin{twierdzenie}[Lemat Fatou]
\label{thm:fatou}
Niech $f_n:X\to[0,\infty]$ będą funkcjami mierzalnymi.
Wtedy
\begin{equation}\label{eq:fatou}
\int_X\liminf_{n\to\infty}f_n(x)\,\mathrm d\mu(x)
\le\liminf_{n\to\infty}\int_Xf_n(x)\,\mathrm d\mu(x).
\end{equation}
Nie wymagamy monotoniczności ciągu ani skończoności miary.
\end{twierdzenie}
\noindent\emph{Dowód:} zob. \atlasref{ME-08::thm:fatou}{Teoria miary, rozdział 8 (dowód tego twierdzenia)}.

\subsection{Twierdzenie Lebesgue'a o zbieżności zmajoryzowanej}


\begin{twierdzenie}[O zbieżności zmajoryzowanej; Lebesgue]
\label{thm:dct}
Niech $f_n:X\to\mathbb R$ będą mierzalne i
$f_n(x)\to f(x)$ dla $\mu$-prawie wszystkich $x$, gdzie
$f:X\to\mathbb R$ jest mierzalna. Załóżmy, że
istnieje nieujemna całkowalna funkcja mierzalna
$g:X\to[0,\infty]$ taka, że
$|f_n(x)|\le g(x)$ dla wszystkich $n$ i
$\mu$-prawie wszystkich $x$, przy czym zbiór
wyjątkowy można wybrać wspólny dla całego ciągu.
Wówczas $f_n$ i $f$ są całkowalne oraz
\begin{equation}\label{eq:dct}
\lim_{n\to\infty}\int_Xf_n(x)\,\mathrm d\mu(x)
=\int_Xf(x)\,\mathrm d\mu(x).
\end{equation}
\end{twierdzenie}
\noindent\emph{Dowód:} zob. \atlasref{ME-08::thm:dct}{Teoria miary, rozdział 8 (dowód tego twierdzenia)}.

\begin{wniosek}[Zbieżność zmajoryzowana w $L^1$]
\label{cor:dct-l1}
Przy założeniach twierdzenia \ref{thm:dct} zachodzi
\[
\lim_{n\to\infty}\int_X|f_n(x)-f(x)|\,\mathrm d\mu(x)=0.
\]
Dla funkcji zespolonych wniosek i twierdzenie
\ref{thm:dct} także zachodzą, jeśli zastąpić
wartość bezwzględną modułem.
\end{wniosek}
\noindent\emph{Dowód:} zob. \atlasref{ME-08::cor:dct-l1}{Teoria miary, rozdział 8 (dowód tego twierdzenia)}.

\begin{wniosek}[Twierdzenie o zbieżności ograniczonej]
\label{cor:bct}
Niech $\mu(X)<\infty$ oraz $f_n,f:X\to\mathbb R$
będą mierzalne. Jeśli $f_n\to f$ prawie wszędzie
i istnieje $M\in[0,\infty)$ takie, że $|f_n|\le M$
prawie wszędzie dla wszystkich $n$, to wszystkie $f_n$ oraz $f$ są całkowalne i
$\int_X f_n\,\mathrm d\mu\to\int_X f\,\mathrm d\mu$ oraz $\int_X|f_n-f|\,\mathrm d\mu\to0$.
\end{wniosek}
\noindent\emph{Dowód:} zob. \atlasref{ME-08::cor:bct}{Teoria miary, rozdział 8 (dowód tego twierdzenia)}.

\subsection{Szeregi funkcji i całek}


\begin{twierdzenie}[Szeregi nieujemnych funkcji mierzalnych]
\label{thm:sum-nonnegative-integrals}
Niech $f_n:X\to[0,\infty]$ będą mierzalne. Wtedy
\begin{equation}\label{eq:sum-nonnegative-integrals}
\int_X\sum_{n=1}^{\infty}f_n(x)\,\mathrm d\mu(x)
=\sum_{n=1}^{\infty}\int_Xf_n(x)\,\mathrm d\mu(x),
\end{equation}
gdzie obie strony mogą być nieskończone.
\end{twierdzenie}
\noindent\emph{Dowód:} zob. \atlasref{ME-08::thm:sum-nonnegative-integrals}{Teoria miary, rozdział 8 (dowód tego twierdzenia)}.

\begin{twierdzenie}[Całkowanie szeregów bezwzględnie sumowalnych]
\label{thm:sum-absolutely-integrable}
Niech $f_n:X\to\mathbb C$ będą mierzalne i
całkowalne oraz niech
$\sum_{n\ge1}\int_X|f_n(x)|\,\mathrm d\mu(x)<\infty$.
Wówczas $\sum_nf_n(x)$ jest bezwzględnie zbieżny
$\mu$-prawie wszędzie. Jeśli $F$ oznacza jego sumę
na zbiorze zbieżności i zero na pozostałym zbiorze,
to $F$ jest całkowalna, a dla sum częściowych $S_N$
zachodzi oszacowanie ogona
\[
 \int_X|F-S_N|\,\mathrm d\mu\le
 \sum_{n>N}\int_X|f_n|\,\mathrm d\mu\longrightarrow0.
\]
Ponadto szereg liczbowy $\sum_n\int_X f_n\,\mathrm d\mu$
jest bezwzględnie zbieżny.
W szczególności
\begin{equation}\label{eq:interchange-series}
\int_XF(x)\,\mathrm d\mu(x)
=\sum_{n=1}^{\infty}\int_Xf_n(x)\,\mathrm d\mu(x).
\end{equation}
\end{twierdzenie}
\begin{proof}
Zob. \atlasref{ME-08::thm:sum-absolutely-integrable}{Teoria miary, rozdział 8},
gdzie wykazano również oszacowanie ogona
i bezwzględną zbieżność szeregu całek.
\end{proof}


\subsection{Którego twierdzenia użyć?}
Zbieżność monotoniczna pozwala obliczać całkę funkcji
nieujemnej za pomocą rosnących obcięć i dopuszcza
nieskończoną wartość całki. Lemat Fatou daje jednostronne
oszacowanie bez monotoniczności. Zbieżność zmajoryzowana
pozwala zamienić granicę z całką i uzyskać zbieżność w $L^1$,
jeżeli istnieje wspólna całkowalna majoranta.
Na przestrzeni o skończonej mierze wystarczy wspólne
ograniczenie stałą. W przypadku szeregu stosujemy
zbieżność monotoniczną do nieujemnych sum częściowych,
a przy funkcjach ze znakiem sprawdzamy bezwzględną
sumowalność całek modułów.

\begin{przyklad}
Na $(0,1)$ z miarą Lebesgue'a $\lambda_1$ niech
$f_n(x)=\min(n,x^{-1/2})$. Funkcje są mierzalne,
nieujemne i $f_n\uparrow x^{-1/2}$. Rozdzielając
przedział w punkcie $n^{-2}$, otrzymujemy
\[
 \int_{(0,1)} f_n\,\mathrm d\lambda_1
 = n\cdot n^{-2}+\int_{n^{-2}}^1 x^{-1/2}\,\mathrm dx
 =\frac1n+2(1-\tfrac1n)=2-\frac1n.
\]
Drugi składnik obliczamy w sensie Riemanna na zwartym
przedziale, korzystając ze
\atlasref{ME-08::thm:riemann-to-lebesgue}{zgodności całek Riemanna i Lebesgue'a (ME-08)}.
Dla $n=1$ drugi składnik wynosi zero.
Twierdzenie \ref{thm:mct} daje więc
$\int_{(0,1)}x^{-1/2}\,\mathrm d\lambda_1(x)=2$.
Jeżeli natomiast $g_n(x)=x^n$ na $(0,1)$, to $g_n\to0$
punktowo i $0\le g_n\le1$. Stała majoranta $1$ jest
całkowalna, ponieważ $\lambda_1((0,1))=1$.
Z twierdzenia \ref{thm:dct} całki dążą zatem do zera.
Sprawdzamy to bezpośrednio: po ciągłym przedłużeniu
$g_n$ na $[0,1]$ korzystamy ze zgodności całek
Riemanna i Lebesgue'a (końce mają miarę zero), skąd
\[
 \int_{(0,1)}x^n\,\mathrm d\lambda_1(x)
 =\left[\frac{x^{n+1}}{n+1}\right]_0^1
 =\frac1{n+1}\longrightarrow0.
\]
\end{przyklad}

Dalsze warianty twierdzeń granicznych, w tym jednostajną
całkowalność i twierdzenie Vitaliego, przedstawiono w
\atlasref{ME-08::thm:vitali-l1}{ME-08: twierdzenie Vitaliego}.
